% =========================================================
%  ppt/04-fisika-gravitasi-lompat-pantul.tex
%  Materi 4 — Fisika Sederhana: Gravitasi, Lompatan, dan Pantulan
%  Bab VI buku ajar
% =========================================================
\documentclass[aspectratio=169,10pt]{beamer}
% =========================================================
%  ppt/config/packages.tex
%  Paket bersama untuk seluruh presentasi kuliah
%  Pengantar Pemrograman Game
% =========================================================

\usepackage[T1]{fontenc}
\usepackage[utf8]{inputenc}
\usepackage[indonesian]{babel}
\usepackage{lmodern}
\usepackage{textcomp}
\usepackage{amsmath}
\usepackage{amssymb}
\usepackage{graphicx}
\usepackage{booktabs}
\usepackage{tabularx}
\usepackage{array}
\usepackage{multicol}
\usepackage{listings}
\usepackage{upquote}
\usepackage{adjustbox}
\usepackage{tcolorbox}
\usepackage{tikz}
\usetikzlibrary{
  arrows.meta,
  positioning,
  shapes.geometric,
  shapes.misc,
  calc,
  fit,
  backgrounds,
  matrix,
  chains,
  decorations.pathreplacing,
  patterns
}

\tcbuselibrary{skins}

% Jangan memuat geometry: ukuran halaman diatur oleh beamer.

% =========================================================
%  ppt/config/theme.tex
%  Satu tema Beamer untuk keempat belas presentasi
%  Pengantar Pemrograman Game
% =========================================================

% --------------------------------------------------------- warna
\definecolor{ubNavy}{HTML}{1F3864}
\definecolor{ubAccent}{HTML}{2E75B6}
\definecolor{ubHead}{HTML}{D9E2F3}
\definecolor{ubSoft}{HTML}{F2F6FC}
\definecolor{ubLine}{HTML}{BFBFBF}
\definecolor{ubText}{HTML}{1F2937}
\definecolor{ubMuted}{HTML}{5A6678}
\definecolor{ubGreen}{HTML}{1E7B45}
\definecolor{ubGreenBg}{HTML}{F1F8F2}
\definecolor{ubRed}{HTML}{B3261E}
\definecolor{ubRedBg}{HTML}{FDF2F1}
\definecolor{ubWarn}{HTML}{A35A05}
\definecolor{ubWarnBg}{HTML}{FFF8EC}
\definecolor{ubGray}{HTML}{F6F6F6}
\definecolor{ubCodeBg}{HTML}{10233F}
\definecolor{ubCodeText}{HTML}{E8EEF7}
\definecolor{ubCodeKey}{HTML}{8FBCF0}
\definecolor{ubCodeCom}{HTML}{8FB37A}
\definecolor{ubCodeStr}{HTML}{D9A46B}
\definecolor{ubCodeNum}{HTML}{7E99BE}

% Nama warna yang dipakai gambar TikZ buku ajar.
\definecolor{ubPrimary}{HTML}{1F3864}
\definecolor{ubAlt}{HTML}{F2F6FC}

% --------------------------------------------------------- dasar
\usetheme{default}
\usefonttheme{professionalfonts}
\setbeamertemplate{navigation symbols}{}
\setbeamerfont{frametitle}{size=\Large,series=\bfseries}
\setbeamerfont{framesubtitle}{size=\small}
\setbeamerfont{footline}{size=\tiny}
\setbeamerfont{block title}{size=\small,series=\bfseries}
\setbeamerfont{block body}{size=\small}
\setbeamerfont{itemize/enumerate body}{size=\small}
\setbeamerfont{itemize/enumerate subbody}{size=\footnotesize}

\setbeamersize{text margin left=0.60cm,text margin right=0.60cm}
\uchyph=0

\setbeamercolor{normal text}{fg=ubText,bg=white}
\setbeamercolor{structure}{fg=ubAccent}
\setbeamercolor{frametitle}{fg=ubNavy,bg=white}
\setbeamercolor{title}{fg=ubNavy}
\setbeamercolor{footline}{fg=ubMuted,bg=white}
\setbeamercolor{itemize item}{fg=ubAccent}
\setbeamercolor{itemize subitem}{fg=ubAccent}
\setbeamercolor{block title}{fg=white,bg=ubNavy}
\setbeamercolor{block body}{fg=ubText,bg=ubSoft}
\setbeamercolor{block title alerted}{fg=white,bg=ubWarn}
\setbeamercolor{block body alerted}{fg=ubText,bg=ubWarnBg}
\setbeamercolor{block title example}{fg=white,bg=ubGreen}
\setbeamercolor{block body example}{fg=ubText,bg=ubGreenBg}
\setbeamercolor{caption name}{fg=ubMuted}

\setbeamertemplate{blocks}[rounded][shadow=false]
\setbeamertemplate{itemize item}{\raisebox{0.10em}{\color{ubAccent}\scriptsize$\blacksquare$}}
\setbeamertemplate{itemize subitem}{\raisebox{0.06em}{\color{ubAccent}\tiny$\blacktriangleright$}}
\setbeamertemplate{enumerate item}{%
  \makebox[1.3em][l]{\color{ubAccent}\insertenumlabel.}}

\addtobeamertemplate{itemize/enumerate body begin}{%
  \setlength{\itemsep}{0.42em}\setlength{\parskip}{0pt}}{}

% --------------------------------------------------------- judul frame
\setbeamertemplate{frametitle}[default][left]

% --------------------------------------------------------- kaki
\newcommand{\materino}{00}
\newcommand{\materifoot}{Pengantar Pemrograman Game \textperiodcentered\ Materi \materino}

\setbeamertemplate{footline}{%
  \begin{beamercolorbox}[wd=\paperwidth,ht=0.015cm,dp=0cm]{footline}%
    \color{ubHead}\rule{\paperwidth}{0.4pt}%
  \end{beamercolorbox}%
  \begin{beamercolorbox}[wd=\paperwidth,ht=0.40cm,dp=0.14cm,%
                         leftskip=0.60cm,rightskip=0.60cm]{footline}%
    \usebeamerfont{footline}\color{ubMuted}\materifoot\hfill\insertframenumber%
  \end{beamercolorbox}%
}

% --------------------------------------------------------- makro buku ajar
\providecommand{\code}[1]{\texttt{\detokenize{#1}}}
\providecommand{\istilah}[1]{\emph{#1}}
\providecommand{\CPMK}[1]{CPMK\,#1}
\providecommand{\CPL}[1]{CPL\,#1}
\providecommand{\SUBCPMK}[1]{Sub-CPMK\,#1}

% --------------------------------------------------------- gaya kode
\lstdefinelanguage{CSharp}{
  morekeywords={abstract,as,base,bool,break,byte,case,catch,char,class,const,
    continue,decimal,default,delegate,do,double,else,enum,event,explicit,
    extern,false,finally,fixed,float,for,foreach,goto,if,implicit,in,int,
    interface,internal,is,lock,long,namespace,new,null,object,operator,out,
    override,params,private,protected,public,readonly,ref,return,sbyte,
    sealed,short,sizeof,stackalloc,static,string,struct,switch,this,throw,
    true,try,typeof,uint,ulong,unchecked,unsafe,ushort,using,virtual,void,
    volatile,while,var,get,set},
  sensitive=true,
  morecomment=[l]{//},
  morecomment=[s]{/*}{*/},
  morestring=[b]",
}

\lstdefinestyle{kode}{
  basicstyle=\ttfamily\footnotesize\color{ubCodeText},
  keywordstyle=\color{ubCodeKey}\bfseries,
  commentstyle=\color{ubCodeCom}\itshape,
  stringstyle=\color{ubCodeStr},
  numberstyle=\tiny\color{ubCodeNum},
  numbers=left,
  numbersep=7pt,
  backgroundcolor=\color{ubCodeBg},
  frame=none,
  showstringspaces=false,
  breaklines=false,
  tabsize=4,
  columns=fullflexible,
  keepspaces=true,
  xleftmargin=12pt,
  xrightmargin=2pt,
  aboveskip=0pt,
  belowskip=0pt,
  upquote=true,
}
\lstset{style=kode}

\newcommand{\kodejelas}[1]{%
  \vspace{0.16cm}%
  {\footnotesize\color{ubMuted}#1\par}%
}
\newcommand{\sumber}[1]{%
  \vspace{0.08cm}%
  {\scriptsize\color{ubMuted}\itshape #1\par}%
}
\newcommand{\awalan}[1]{%
  {\small\color{ubMuted}#1\par}\vspace{0.18cm}%
}

% --------------------------------------------------------- kotak
\newtcolorbox{kotak}[1][]{enhanced, colback=ubSoft, colframe=ubAccent,
  boxrule=0pt, leftrule=2.2pt, arc=1pt, left=6pt, right=6pt, top=4pt,
  bottom=4pt, fontupper=\footnotesize, #1}
\newtcolorbox{peringatan}[1][]{enhanced, colback=ubWarnBg, colframe=ubWarn,
  boxrule=0pt, leftrule=2.2pt, arc=1pt, left=6pt, right=6pt, top=4pt,
  bottom=4pt, fontupper=\footnotesize, #1}
\newtcolorbox{kiat}[1][]{enhanced, colback=ubGreenBg, colframe=ubGreen,
  boxrule=0pt, leftrule=2.2pt, arc=1pt, left=6pt, right=6pt, top=4pt,
  bottom=4pt, fontupper=\footnotesize, #1}

% --------------------------------------------------------- sampul
% \deckcover{nomor}{minggu}{judul}{sub-cpmk}{teknologi}{sumber}{blok}
\newcommand{\deckcover}[7]{%
\begin{frame}[plain]
\begin{tikzpicture}[remember picture, overlay]
  \fill[ubNavy] (current page.south west) rectangle ([xshift=6.00cm]current page.north west);
  \fill[ubAccent] ([xshift=6.00cm]current page.south west) rectangle ([xshift=6.08cm]current page.north west);

  \node[anchor=north west, text width=4.72cm, align=left, inner sep=0pt]
    at ($(current page.north west)+(0.66cm,-0.60cm)$) {%
      {\color{white}\bfseries\fontsize{10.2}{12.2}\selectfont
       Universitas Bale Bandung}\\[4pt]
      {\color{ubHead}\fontsize{7.6}{9.4}\selectfont
       Fakultas Teknologi Informasi}\\[1pt]
      {\color{ubHead}\fontsize{7.6}{9.4}\selectfont
       Program Studi Teknik Informatika}\\[7pt]
      {\color{ubAccent}\rule{4.0cm}{0.5pt}}%
    };

  \node[anchor=north west, text width=4.72cm, align=left, inner sep=0pt]
    at ($(current page.north west)+(0.66cm,-2.72cm)$) {%
      {\color{ubAccent}\fontsize{6.4}{7.8}\selectfont MATA KULIAH}\\[3pt]
      {\color{white}\bfseries\fontsize{9.8}{12.0}\selectfont
       Pengantar Pemrograman Game}\\[3pt]
      {\color{ubHead}\fontsize{7.6}{9.4}\selectfont
       MWCA009 \textperiodcentered\ 3 SKS}\\[8pt]
      {\color{ubHead}\fontsize{7.6}{9.4}\selectfont
       Semester Ganjil 2026/2027}%
    };

  \node[anchor=south west, text width=4.72cm, align=left, inner sep=0pt]
    at ($(current page.south west)+(0.66cm,0.52cm)$) {%
      {\color{ubAccent}\fontsize{6.4}{7.8}\selectfont DOSEN PENGAMPU}\\[3pt]
      {\color{white}\fontsize{8.4}{10.4}\selectfont
       Cecep Suwanda, S.Si., M.Kom.}%
    };

  \node[anchor=north west, text width=8.95cm, align=left, inner sep=0pt]
    at ($(current page.north east)+(-9.40cm,-0.92cm)$) {%
      {\color{ubAccent}\bfseries\fontsize{8.2}{10.2}\selectfont
       MATERI #1\ \ \textperiodcentered\ \ MINGGU #2}\\[8pt]
      {\color{ubNavy}\bfseries\fontsize{19.4}{23.4}\selectfont
       \hyphenpenalty=10000 \exhyphenpenalty=10000 \raggedright #3}%
    };

  \node[anchor=north west, text width=8.95cm, align=left, inner sep=0pt]
    at ($(current page.north east)+(-9.40cm,-4.30cm)$) {%
      \renewcommand{\arraystretch}{1.30}%
      \begin{tabular}{@{}>{\color{ubMuted}\fontsize{7.3}{9.0}\selectfont}p{1.85cm}
                        >{\color{ubText}\fontsize{8.1}{10.0}\selectfont}p{6.7cm}@{}}
        Sub-CPMK      & #4 \\
        Teknologi     & #5 \\
        Sumber materi & #6 \\
        Blok          & #7 \\
      \end{tabular}%
    };
\end{tikzpicture}
\end{frame}%
}

% --------------------------------------------------------- pembatas bagian
\newcommand{\sectionframe}[3]{%
\begin{frame}[plain]
\begin{tikzpicture}[remember picture, overlay]
  \fill[ubNavy] (current page.south west) rectangle (current page.north east);
  \fill[ubAccent] (current page.north west) rectangle ([xshift=0.10cm]current page.south west);
  \node[anchor=west, inner sep=0pt]
    at ($(current page.west)+(1.00cm,1.45cm)$)
    {\color{ubNavy!52!white}\bfseries\fontsize{58}{58}\selectfont #1};
  \node[anchor=west, inner sep=0pt, text width=11.6cm, align=left]
    at ($(current page.west)+(1.05cm,-0.55cm)$) {%
      {\color{ubAccent}\bfseries\fontsize{8.0}{10.0}\selectfont BAGIAN}\\[3pt]
      {\color{white}\bfseries\fontsize{20.5}{24.5}\selectfont #2}\\[5pt]
      {\color{ubHead}\small #3}%
    };
  \node[anchor=south east, inner sep=0pt]
    at ($(current page.south east)+(-0.60cm,0.32cm)$)
    {\color{ubHead}\fontsize{7.4}{9}\selectfont #1 / 6
     \ \ \textperiodcentered\ \ \insertframenumber};
\end{tikzpicture}
\end{frame}%
}

% --------------------------------------------------------- peta pertemuan
\newcommand{\blokpeta}[6]{%
  \ifnum#2=#3
    \node[rounded corners=3pt, fill=ubNavy, draw=ubNavy, line width=0.5pt,
          text width=3.02cm, minimum height=1.62cm, align=left, inner sep=5pt]
      at (#1,0) {%
        {\color{ubAccent}\bfseries\fontsize{6.8}{8}\selectfont #4}\\[1pt]
        {\color{white}\bfseries\fontsize{7.6}{9}\selectfont #5}\\[2pt]
        {\color{ubHead}\fontsize{6.4}{7.6}\selectfont #6}};
  \else
    \node[rounded corners=3pt, fill=ubSoft, draw=ubHead, line width=0.5pt,
          text width=3.02cm, minimum height=1.62cm, align=left, inner sep=5pt]
      at (#1,0) {%
        {\color{ubAccent}\bfseries\fontsize{6.8}{8}\selectfont #4}\\[1pt]
        {\color{ubNavy}\bfseries\fontsize{7.6}{9}\selectfont #5}\\[2pt]
        {\color{ubMuted}\fontsize{6.4}{7.6}\selectfont #6}};
  \fi
}

% \petapertemuan{blok aktif}{sebelumnya}{sesudahnya}{keluaran}
% \petapertemuan{blok aktif}{sebelumnya}{pertemuan ini}{sesudahnya}{keluaran}
\newcommand{\petapertemuan}[5]{%
\begin{frame}{Peta Pertemuan}
{\scriptsize\color{ubMuted}Kedudukan materi dalam empat blok perkuliahan}\par
\vspace{0.10cm}
\centering
\begin{tikzpicture}
  \blokpeta{-5.58}{1}{#1}{BLOK I}{Materi 1--5}{Game loop, vektor,\\fisika, collision}
  \blokpeta{-1.86}{2}{#1}{BLOK II}{Materi 6--9}{AI NPC: chase, FSM,\\pathfinding A*}
  \blokpeta{1.86}{3}{#1}{BLOK III}{Materi 10--11}{Unity: scene,\\level, physics}
  \blokpeta{5.58}{4}{#1}{BLOK IV}{Materi 12--14}{Aset visual, audio,\\UI, playtest}
\end{tikzpicture}
\par\vspace{0.30cm}
\begin{columns}[T,onlytextwidth]
  \begin{column}{0.315\textwidth}
    \begin{beamercolorbox}[wd=\linewidth,sep=4pt,rounded=true]{block body}
      {\color{ubMuted}\fontsize{6.6}{8}\selectfont SEBELUMNYA}\\[2pt]
      {\footnotesize #2}
    \end{beamercolorbox}
  \end{column}
  \begin{column}{0.315\textwidth}
    \begin{beamercolorbox}[wd=\linewidth,sep=4pt,rounded=true]{block title}
      {\color{ubHead}\fontsize{6.6}{8}\selectfont PERTEMUAN INI}\\[2pt]
      {\footnotesize #3}
    \end{beamercolorbox}
  \end{column}
  \begin{column}{0.315\textwidth}
    \begin{beamercolorbox}[wd=\linewidth,sep=4pt,rounded=true]{block body}
      {\color{ubMuted}\fontsize{6.6}{8}\selectfont SESUDAHNYA}\\[2pt]
      {\footnotesize #4}
    \end{beamercolorbox}
  \end{column}
\end{columns}
\par\vspace{0.26cm}
{\footnotesize\color{ubText}\textbf{Keluaran pertemuan.} #5\par}
\end{frame}%
}

% --------------------------------------------------------- penutup
\newcommand{\penutup}[2]{%
\begin{frame}[plain]
\begin{tikzpicture}[remember picture, overlay]
  \fill[ubNavy] (current page.south west) rectangle (current page.north east);
  \fill[ubAccent] (current page.north west) rectangle ([xshift=0.10cm]current page.south west);

  \node[anchor=north west, text width=13.6cm, align=left, inner sep=0pt]
    at ($(current page.north west)+(1.00cm,-0.95cm)$) {%
      {\color{ubAccent}\bfseries\fontsize{8.2}{10}\selectfont ASESMEN MINGGU INI}\\[6pt]
      {\color{white}\bfseries\fontsize{24}{28}\selectfont #1}%
    };

  \node[anchor=north west, text width=13.6cm, align=left, inner sep=0pt]
    at ($(current page.north west)+(1.00cm,-4.55cm)$) {%
      {\color{ubNavy!0!ubHead}\rule{13.4cm}{0.4pt}}\\[7pt]
      {\color{ubAccent}\bfseries\fontsize{7.4}{9}\selectfont
       PENGAIT KE MATERI BERIKUTNYA}\\[4pt]
      {\color{ubHead}\fontsize{10.6}{13.2}\selectfont #2}%
    };

  \node[anchor=south east, inner sep=0pt]
    at ($(current page.south east)+(-0.60cm,0.32cm)$)
    {\color{ubHead}\fontsize{7.4}{9}\selectfont \insertframenumber};
\end{tikzpicture}
\end{frame}%
}



\renewcommand{\materino}{04}

\begin{document}

\deckcover{04}{4}
  {Fisika Sederhana: Gravitasi, Lompatan, dan Pantulan}
  {1.3}
  {Python / Pygame}
  {Bab VI buku ajar}
  {Blok I \textperiodcentered\ Materi 1--5}

% --------------------------------------------------------- capaian
\begin{frame}{Capaian Pertemuan}
\setbeamerfont{itemize/enumerate body}{size=\scriptsize}
\awalan{Sub-CPMK dan indikator dari RPS, serta tujuan Bab VI buku ajar.}
\begin{columns}[T,onlytextwidth]
  \begin{column}{0.52\textwidth}
    \begin{beamercolorbox}[wd=\linewidth,sep=5pt,rounded=true]{block title}
      {\fontsize{6.8}{8}\selectfont\color{ubHead}SUB-CPMK 1.3}\\[2pt]
      {\footnotesize Menerapkan logika fisika sederhana (\emph{gravitasi},
       lompatan, pantulan) dalam simulasi 2D.}
    \end{beamercolorbox}
    \vspace{0.16cm}

    {\fontsize{7}{8.5}\selectfont\bfseries\color{ubAccent}INDIKATOR}\par
    \vspace{0.06cm}
    \begin{itemize}
      \item Menerapkan \emph{integrasi Euler} untuk memperbarui kecepatan dan
            posisi.
      \item Mengimplementasikan lompatan yang hanya aktif ketika objek berada
            di tanah.
      \item Menghasilkan pantulan dengan \emph{koefisien restitusi} sederhana.
    \end{itemize}
  \end{column}
  \begin{column}{0.45\textwidth}
    {\fontsize{7}{8.5}\selectfont\bfseries\color{ubAccent}TUJUAN BAB VI}\par
    \vspace{0.08cm}
    \begin{itemize}
      \item Menerapkan integrasi Euler dua tingkat untuk memperbarui kecepatan
            lalu posisi.
      \item Menjelaskan mengapa arah percepatan gravitasi bernilai positif
            pada koordinat layar.
      \item Menghentikan benda di lantai, dan menjelaskan mengapa pemain tidak
            boleh memantul ketika mendarat.
      \item Mengimplementasikan lompatan yang hanya bekerja ketika benda berada
            di tanah.
      \item Memilih nilai gravitasi dan kecepatan lompatan untuk memperoleh
            tinggi lompatan yang dikehendaki.
      \item Menghasilkan pantulan dengan koefisien restitusi, dan menjelaskan
            mengapa tinggi pantulan menyusut sebesar kuadrat koefisien itu.
    \end{itemize}
  \end{column}
\end{columns}
\end{frame}

% --------------------------------------------------------- peta
\petapertemuan{1}
  {Materi 3 \textperiodcentered\ Gerak Objek, \texttt{Vector2}, Input, dan
   \emph{Delta Time}}
  {Materi 4 \textperiodcentered\ Fisika Sederhana: Gravitasi, Lompatan, dan
   Pantulan}
  {Materi 5 \textperiodcentered\ \emph{Sprite}, \emph{Group}, dan
   \emph{Collision Detection}}
  {Prototipe permainan yang bendanya jatuh, melompat, dan memantul; tinggi
   lompatan terukur sesuai dengan hitungan.}

% ========================================================= bagian 1
\section{Percepatan sebagai Perubahan Kecepatan}
\sectionframe{01}{Percepatan sebagai\\Perubahan Kecepatan}
  {Tingkat kedua pada tahap pembaruan keadaan}

\begin{frame}{Tingkat Kedua pada Pembaruan Keadaan}
\awalan{Pada bab sebelumnya, kecepatan benda berubah hanya karena tombol ditekan
atau dilepas. Bab ini membahas perubahan kecepatan yang terjadi dengan
sendirinya, yaitu karena adanya percepatan.}
\begin{itemize}
  \item Penerapannya hanya menambah satu baris dibandingkan Bab V: kecepatan
        tegak bertambah setiap bingkai, dan penambahannya sebanding dengan
        selang waktu antar bingkai (\emph{delta time}).
  \item Kedua baris dikerjakan berurutan; baris kedua memakai kecepatan yang
        baru saja diperbarui pada baris pertama. Itulah sebabnya benda yang
        jatuh bergerak makin cepat.
\end{itemize}
\[
  v \leftarrow v + a\,\Delta t
  \qquad\qquad
  p \leftarrow p + v\,\Delta t
\]
\begin{kotak}
Urutan itu bukan pilihan gaya penulisan, melainkan kaidah \emph{integrasi
Euler} semi-implisit. Apabila urutannya dibalik, benda tampak jatuh lebih cepat
daripada seharusnya, dan kesalahan itu menumpuk ke satu arah sehingga tidak
lagi terlihat sebagai kesalahan kecil.
\end{kotak}
\end{frame}

\begin{frame}{Arah Gravitasi pada Koordinat Layar}
\awalan{Pada koordinat layar, sumbu tegak menuju ke bawah. Karena itu,
percepatan \emph{gravitasi} yang menarik benda ke bawah dinyatakan sebagai
nilai positif.}
\[
  a_y = +g > 0, \qquad\qquad v_{\text{lompat}} < 0
\]
\begin{itemize}
  \item Kecepatan lompatan bernilai negatif karena mengarah ke atas. Kedua
        tanda itu mudah tertukar, dan gejalanya khas.
  \item Benda yang seharusnya jatuh justru melayang ke atas, atau lompatan yang
        seharusnya ke atas justru menekan benda ke lantai.
  \item Nilai gravitasi diukur dalam piksel per detik kuadrat, bukan meter per
        detik kuadrat; nilai yang sering dipakai berkisar 1000--2500 piksel per
        detik kuadrat.
\end{itemize}
\begin{kotak}
Pilihan nilainya bergantung pada skala permainan dan pada kecepatan lompatan,
sebagaimana dibahas pada bagian tentang lompatan.
\end{kotak}
\end{frame}

\begin{frame}[fragile]{Satu Fungsi untuk Semua Benda}
\lstinputlisting[language=Python,firstline=69,lastline=81]
  {../buku_ajar/code/python/ch06_fisika.py}
\kodejelas{Dua tingkat terlihat berurutan: percepatan mengubah kecepatan, lalu
kecepatan mengubah posisi. Gravitasi dan restitusi diberikan sebagai argumen,
bukan ditulis di dalam fungsi, sehingga satu fungsi melayani semua benda dan
perbaikan cukup dilakukan di satu tempat.}
\end{frame}

% ========================================================= bagian 2
\section{Gravitasi dan Batas Lantai}
\sectionframe{02}{Gravitasi dan Batas Lantai}
  {Lantai sebagai batas, dan persoalan ketika benda menyentuhnya}

\begin{frame}{Lantai sebagai Batas Bawah}
\awalan{Benda yang jatuh perlu berhenti di suatu tempat. Lantai dinyatakan
sebagai satu nilai tetap: jarak tepi atasnya dari tepi atas jendela.}
\begin{itemize}
  \item Pada program bab ini lantai berada 60 piksel di atas tepi bawah jendela,
        menyisakan ruang untuk menampilkan keterangan.
  \item Setiap bingkai, sesudah posisi diperbarui, program memeriksa apakah
        benda sudah menyentuh atau melewati lantai.
  \item Pemeriksaan memakai tanda ketidaksamaan, bukan kesamaan. Benda yang
        bergerak cepat melampaui lantai dalam satu bingkai, sehingga posisinya
        tidak pernah tepat sama dengan nilai lantai.
  \item Posisi dikembalikan ke permukaan lantai sebelum kecepatan diubah, agar
        benda tidak terbenam semakin dalam.
  \item Kecepatan jatuh dibatasi pada nilai tertentu, misalnya 1200 piksel per
        detik, untuk mencegah benda menembus lantai. Nilai batas itu disebut
        kecepatan akhir (\emph{terminal velocity}).
\end{itemize}
\end{frame}

\begin{frame}{Mengapa Pemain Tidak Boleh Memantul}
\awalan{Pemain dan bola memakai fungsi pembaruan yang sama, tetapi satu nilai
argumennya berbeda. Perhatikan akibatnya.}
{\scriptsize
\renewcommand{\arraystretch}{1.14}
\begin{tabularx}{\textwidth}{@{}>{\raggedright\arraybackslash}p{1.8cm}
    >{\centering\arraybackslash}p{1.7cm}
    >{\raggedright\arraybackslash}p{4.5cm}
    >{\raggedright\arraybackslash}X@{}}
\toprule
\textbf{Benda} & \textbf{Restitusi} & \textbf{Perilaku saat mendarat}
  & \textbf{Alasan} \\
\midrule
Pemain & $0{,}0$ &
  Berhenti seketika; kecepatan tegaknya menjadi nol. &
  Kendali pemain menuntut ketepatan. Pemain yang memantul tidak dapat
  mengarahkan geraknya, dan permainan terasa tidak dapat diandalkan. \\
\addlinespace
Bola & $0{,}8$ &
  Memantul dengan kecepatan 80 persen dari kecepatan datangnya. &
  Pantulan memang dikehendaki; kehilangan kecepatan membuat bola akhirnya
  berhenti dengan sendirinya. \\
\bottomrule
\end{tabularx}}
\vspace{0.14cm}
\begin{peringatan}
Kesalahan yang sering terjadi adalah memberi koefisien restitusi kepada pemain
``supaya terlihat hidup''. Akibatnya pemain memantul setiap kali mendarat dan
kendalinya terasa lepas. Untuk kesan hidup, tambahkan percepatan dan
perlambatan pada gerak mendatar, bukan pantulan pada gerak tegak.
\end{peringatan}
\end{frame}

\begin{frame}[fragile]{Mengenali Benda yang Menyentuh Lantai}
\lstinputlisting[language=Python,firstline=86,lastline=94]
  {../buku_ajar/code/python/ch06_fisika.py}
\kodejelas{Pemeriksaan memakai \texttt{>=} agar benda cepat tetap terdeteksi;
posisi dikembalikan sebelum kecepatan diubah; syarat \texttt{velocity.y > 0}
memastikan hanya benda yang bergerak ke bawah yang dipantulkan; dan
\texttt{MIN\_BOUNCE\_SPEED} menghentikan pantulan yang terlalu kecil. Keadaan
menyentuh lantai diperiksa dengan dua syarat: badan bersandar pada lantai
dengan toleransi setengah piksel, dan kecepatan tegak tidak lagi mengarah ke
bawah.}
\end{frame}

% ========================================================= bagian 3
\section{Lompatan sebagai Impuls}
\sectionframe{03}{Lompatan sebagai Impuls}
  {Perubahan kecepatan seketika, syarat, dan pemilihan nilainya}

\begin{frame}{Perubahan Kecepatan Seketika}
\awalan{Lompatan adalah \emph{jump impulse}: perubahan kecepatan yang terjadi
seketika, bukan gaya yang bekerja selama beberapa waktu.}
\begin{itemize}
  \item Kecepatan lompatan diberikan sebagai nilai negatif karena mengarah ke
        atas, dan nilainya \textbf{diganti}, bukan ditambahkan: satu baris
        \texttt{player\_velocity.y = -jump\_speed}.
  \item Setelah itu gravitasi bekerja seperti biasa: benda melambat naik,
        berhenti di puncak, lalu jatuh kembali.
\end{itemize}
\vspace{0.04cm}
\centering
\adjustbox{max width=0.80\linewidth,max height=0.40\textheight}{%
  % chapters/06/figures/gravitasi-lompatan.tikz
% Lintasan lompatan: kecepatan awal mengarah ke atas, puncak dengan kecepatan
% tegak nol, lalu jatuh dengan kecepatan yang makin besar karena gravitasi.
\begin{tikzpicture}[
  lintasan/.style={very thick, draw=ubPrimary},
  panah/.style={-{Stealth[length=2.4mm]}, very thick, draw=ubAccent},
  panahgrav/.style={-{Stealth[length=2.4mm]}, very thick, draw=ubRed},
  lantai/.style={thick, draw=ubLine},
  label/.style={font=\scriptsize, align=left},
  keterangan/.style={font=\scriptsize\itshape, ubPrimary, align=left}
]

% Lantai.
\draw[lantai] (0,0) -- (8.0,0);

% Lintasan: naik, mencapai puncak, lalu turun.
\draw[lintasan] (0.8,0) .. controls (1.35,1.9) and (1.85,2.6) .. (2.4,2.6);
\draw[lintasan] (2.4,2.6) .. controls (2.95,2.6) and (3.45,1.9) .. (4.0,0);

% Tinggi lompatan.
\draw[ubGreen, dashed] (0.8,0) -- (0.8,2.6);
\draw[ubGreen, dashed] (0.8,2.6) -- (1.55,2.6);
\node[label, ubGreen, anchor=east] at (0.72,1.3)
  {$h = \dfrac{v_0^2}{2g}$};

% Kecepatan awal, mengarah ke atas.
\draw[panah] (0.98,0.3) -- (1.12,1.45);
\node[label, ubAccent, anchor=south east] at (0.95,0.95) {$v_0$};

% Puncak: kecepatan tegak bernilai nol.
\fill[ubPrimary] (2.4,2.6) circle (1.8pt);
\node[label, anchor=south] at (2.4,2.68) {kecepatan tegak nol};

% Kecepatan saat jatuh, mengarah ke bawah dan makin besar.
\draw[panah] (3.88,1.5) -- (3.72,0.35);
\node[label, ubAccent, anchor=north east] at (3.6,1.05) {makin cepat};

% Arah gravitasi yang selalu ke bawah.
\draw[panahgrav] (5.6,2.3) -- (5.6,0.9);
\node[label, ubRed, anchor=west] at (5.75,1.6)
  {gravitasi $g$,\\selalu ke bawah};

% Keterangan arah pada koordinat layar.
\node[keterangan, anchor=north west, text width=52mm] at (5.2,-0.25)
  {Pada koordinat layar, sumbu tegak menuju ke bawah, sehingga gerak ke atas
   bernilai negatif.};
\end{tikzpicture}
}
\par
\sumber{Lintasan lompatan: kecepatan awal $v_0$, puncak dengan kecepatan tegak
nol, lalu jatuh makin cepat karena gravitasi.}
\end{frame}

\begin{frame}[fragile]{Lompatan Dibaca sebagai Peristiwa}
\awalan{Bila lompatan dibaca dengan \texttt{pygame.key.get\_pressed()},
syaratnya selalu benar selama tombol ditahan, sehingga benda melompat
berulang-ulang dan pemain dapat terbang.}
\lstinputlisting[language=Python,firstline=186,lastline=196]
  {../buku_ajar/code/python/ch06_fisika.py}
\kodejelas{Peristiwa \texttt{KEYDOWN} hanya terjadi sekali pada setiap
penekanan. Syarat \texttt{if on\_ground} mencegah lompatan di udara, sedangkan
\texttt{on\_ground = False} segera sesudah lompatan mencegah lompatan kedua
pada bingkai berikutnya.}
\end{frame}

\begin{frame}{Memilih Gravitasi dan Kecepatan Lompatan}
\awalan{Tinggi lompatan ditentukan oleh gravitasi dan kecepatan lompatan.
Hubungannya dapat dihitung, dan dapat dibalik untuk keperluan perancangan.}
\[
  h = \frac{v_0^2}{2g}, \qquad\qquad
  t = \frac{2v_0}{g}, \qquad\qquad
  v_0 = \sqrt{2gh}
\]
{\scriptsize
\renewcommand{\arraystretch}{1.14}
\begin{tabularx}{\textwidth}{@{}>{\centering\arraybackslash}p{1.9cm}
    >{\centering\arraybackslash}p{3.0cm}
    >{\centering\arraybackslash}p{2.2cm}
    >{\raggedright\arraybackslash}X@{}}
\toprule
\textbf{Gravitasi} & \textbf{Kecepatan lompatan} & \textbf{Lama di udara}
  & \textbf{Rasa gerak} \\
\midrule
1200 & 600 piksel/detik & 1,00 detik & Melayang, seperti di bulan. \\
\addlinespace
1800 & 735 piksel/detik & 0,82 detik & Sedang, banyak dipakai permainan
  platform. \\
\addlinespace
2400 & 849 piksel/detik & 0,71 detik & Gesit dan tegas, cocok untuk permainan
  yang cepat. \\
\bottomrule
\end{tabularx}}
\vspace{0.10cm}
\begin{kotak}
Ketiga pasangan itu menghasilkan tinggi 150 piksel, tetapi rasanya berbeda.
Tinggi lompatan saja tidak menentukan rasa gerak; yang menentukan adalah
perbandingan antara gravitasi dan kecepatan lompatan, yaitu lama benda berada
di udara.
\end{kotak}
\end{frame}

% ========================================================= bagian 4
\section{Pantulan dengan Koefisien Restitusi}
\sectionframe{04}{Pantulan dengan Koefisien Restitusi}
  {Membalik arah, mengurangi kecepatan, dan menyusutkan tinggi}

\begin{frame}{Membalik Arah dan Mengurangi Kecepatan}
\awalan{Pantulan terjadi dalam dua langkah: arah gerak dibalik, lalu besarnya
dikurangi. Keduanya dilakukan sekaligus dengan satu perhitungan.}
\[
  v' = -e\,v
\]
\begin{itemize}
  \item Tanda negatif membalik arah, sedangkan pengali $e$ adalah
        \emph{koefisien restitusi} yang nilainya terletak antara nol dan satu.
  \item $e = 0$: tidak ada pantulan sama sekali, benda berhenti seketika.
        Nilai ini dipakai untuk pemain.
  \item $e = 1$: pantulan sempurna, kecepatan tidak berkurang dan benda
        memantul selamanya. Nilai ini jarang dipakai karena membuat permainan
        tidak pernah tenang.
  \item $0 < e < 1$, misalnya $0{,}8$: benda kehilangan sebagian kecepatannya
        setiap pantulan sehingga akhirnya berhenti.
\end{itemize}
\begin{kotak}
Pada program bab ini pemain memakai $e = 0$ dan bola memakai $e = 0{,}8$.
Perbedaan perilaku keduanya hanya berasal dari satu nilai itu.
\end{kotak}
\end{frame}

\begin{frame}{Tinggi Pantulan Menyusut Kuadrat}
\awalan{Kecepatan setelah pantulan adalah $e$ kali kecepatan sebelumnya,
sedangkan tinggi yang dapat dicapai sebanding dengan kuadrat kecepatan.}
\[
  h' = e^2 h
\]
\vspace{0.04cm}
\centering
\adjustbox{max width=0.95\linewidth,max height=0.42\textheight}{%
  % chapters/06/figures/pantulan.tikz
% Pantulan dengan koefisien restitusi 0,8: setiap pantulan menyisakan 64
% persen tinggi sebelumnya, sebab tinggi sebanding dengan kuadrat kecepatan.
\begin{tikzpicture}[
  bola/.style={circle, draw=ubAccent, thick, fill=ubAlt, inner sep=2pt},
  lantai/.style={thick, draw=ubLine},
  label/.style={font=\scriptsize, align=left},
  lintasan/.style={very thick, draw=ubPrimary},
  ukur/.style={draw=ubGreen, dashed, thin}
]

\draw[lantai] (0,0) -- (8.2,0);

% Jatuh dari ketinggian h.
\draw[lintasan] (0,3.2) .. controls (0.75,1.5) and (1.0,0.5) .. (1.2,0);

% Tiga pantulan berikutnya, tinggi menyusut menjadi 0,64 kali setiap kali.
\draw[lintasan] (1.2,0) .. controls (1.75,1.7) and (2.25,2.05) .. (2.6,2.05);
\draw[lintasan] (2.6,2.05) .. controls (2.95,2.05) and (3.5,1.6) .. (4.0,0);

\draw[lintasan] (4.0,0) .. controls (4.35,1.1) and (4.9,1.31) .. (5.1,1.31);
\draw[lintasan] (5.1,1.31) .. controls (5.35,1.31) and (5.75,1.0) .. (6.2,0);

\draw[lintasan] (6.2,0) .. controls (6.4,0.7) and (6.8,0.84) .. (6.9,0.84);
\draw[lintasan] (6.9,0.84) .. controls (7.05,0.84) and (7.3,0.6) .. (7.7,0);

% Bola pada puncak setiap pantulan.
\node[bola] at (0,3.2) {};
\node[bola] at (2.6,2.05) {};
\node[bola] at (5.1,1.31) {};
\node[bola] at (6.9,0.84) {};

% Garis ukur tinggi setiap pantulan.
\draw[ukur] (0,3.2) -- (8.6,3.2);
\draw[ukur] (0,2.05) -- (8.6,2.05);
\draw[ukur] (0,1.31) -- (8.6,1.31);
\draw[ukur] (0,0.84) -- (8.6,0.84);

\node[label, ubGreen, anchor=west] at (8.7,3.2) {$h$};
\node[label, ubGreen, anchor=west] at (8.7,2.05) {$0{,}64\,h$};
\node[label, ubGreen, anchor=west] at (8.7,1.31) {$0{,}41\,h$};
\node[label, ubGreen, anchor=west] at (8.7,0.84) {$0{,}26\,h$};

\node[label, anchor=north west, text width=150mm] at (0,-0.35)
  {Dengan koefisien restitusi $e = 0{,}8$, kecepatan setelah setiap pantulan
   tinggal $0{,}8$ kali, sedangkan tinggi yang dicapai tinggal
   $e^2 = 0{,}64$ kali, sebab tinggi sebanding dengan kuadrat kecepatan.};
\end{tikzpicture}
}
\par
\sumber{Untuk $e = 0{,}8$, tinggi tinggal $0{,}64$ kali setiap pantulan.}
\end{frame}

\begin{frame}{Urutan Penanganan yang Benar}
\awalan{Satu hal yang mudah terlewat adalah urutan pekerjaan ketika benda
menyentuh lantai.}
\begin{enumerate}
  \item Posisi dikembalikan ke permukaan lantai \textbf{lebih dahulu}, baru
        kecepatan diubah. Bila dibalik, benda membalik arah saat masih terbenam:
        tinggi pantulan menyusut lebih cepat dan benda dapat tersangkut.
  \item Syarat \texttt{velocity.y > 0} memastikan hanya benda yang sedang
        bergerak ke bawah yang dipantulkan. Tanpa syarat itu, benda yang baru
        memantul akan dibalik lagi dan terperangkap di lantai.
  \item Pada dinding, yang dibalik adalah komponen mendatar. \texttt{abs()}
        menentukan arah keluar, sebab tanda negatif sederhana akan membalik
        benda yang sudah bergerak menjauhi dinding.
  \item Pantulan yang kecepatannya di bawah 60 piksel per detik dihentikan agar
        benda tidak bergetar di permukaan lantai (\emph{jitter}).
\end{enumerate}
\end{frame}

% ========================================================= bagian 5
\section{Program Fisika Lengkap}
\sectionframe{05}{Program Fisika Lengkap}
  {Dua benda, satu hukum, dan pengukuran tinggi lompatan}

\begin{frame}{Dua Benda, Satu Hukum}
\awalan{Program bab ini menyatukan ketiga unsur fisika menjadi satu program
yang memuat dua benda dengan perilaku berbeda, tetapi dihitung oleh fungsi
yang sama.}
\begin{itemize}
  \item \textbf{Pemain} --- digerakkan tombol mendatar, jatuh karena
        gravitasi, berhenti di lantai karena $e = 0$, dan dapat melompat ketika
        menyentuh lantai. Benda ini memperlihatkan indikator (a) dan (b).
  \item \textbf{Bola} --- bergerak sendiri, memantul di lantai dan di kedua
        dinding dengan $e = 0{,}8$, sehingga pantulannya makin lama makin
        pendek. Benda ini memperlihatkan indikator (c).
  \item Perbedaan perilaku keduanya hanya berasal dari dua hal: siapa yang
        menentukan kecepatan mendatarnya, dan berapa nilai koefisien
        restitusinya.
\end{itemize}
\begin{kotak}
Karena hukum fisika ditulis sebagai satu fungsi, seluruh benda dijamin
mengikuti hukum yang sama. Apabila ditemukan kesalahan pada perhitungan
gravitasi, perbaikannya cukup dilakukan di satu tempat.
\end{kotak}
\end{frame}

\begin{frame}[fragile]{Satu Fungsi, Dua Pemanggilan}
\lstinputlisting[language=Python,firstline=215,lastline=224]
  {../buku_ajar/code/python/ch06_fisika.py}
\kodejelas{Pemanggilan pertama memakai \texttt{PLAYER\_RESTITUTION}, yang kedua
memakai \texttt{args.restitution}; selebihnya identik. Keadaan menyentuh lantai
diperiksa sesudahnya dengan dua syarat: badan bersandar pada lantai dengan
toleransi setengah piksel, dan kecepatan tegak tidak lagi mengarah ke bawah.}
\end{frame}

\begin{frame}{Mengukur Tinggi Lompatan}
\awalan{Tinggi terukur dibandingkan dengan hitungan pada empat laju bingkai.}
{\scriptsize
\renewcommand{\arraystretch}{1.14}
\begin{tabularx}{\textwidth}{@{}>{\centering\arraybackslash}p{2.2cm}
    >{\centering\arraybackslash}p{3.0cm}
    >{\centering\arraybackslash}p{3.4cm}
    >{\raggedright\arraybackslash}X@{}}
\toprule
\textbf{Batas laju} & \textbf{Tinggi terukur} & \textbf{Menurut hitungan}
  & \textbf{Galat} \\
\midrule
30 FPS  & 124,6 piksel & 124,7 piksel & 11,4 piksel (8,4\%) \\
\addlinespace
60 FPS  & 130,3 piksel & 130,3 piksel & 5,8 piksel (4,3\%) \\
\addlinespace
120 FPS & 133,1 piksel & 133,2 piksel & 2,9 piksel (2,1\%) \\
\addlinespace
240 FPS & 134,5 piksel & 134,7 piksel & 1,5 piksel (1,1\%) \\
\bottomrule
\end{tabularx}}
\sumber{Gravitasi 1800 piksel per detik kuadrat dan kecepatan lompatan 700
piksel per detik, sehingga tinggi yang seharusnya 136,1 piksel.}
\begin{kotak}
Hasil pengukuran sesuai dengan hitungan; selisihnya di bawah 0,2 piksel pada
setiap laju. Simulasi selalu sedikit lebih rendah daripada lintasan sebenarnya;
itulah galat integrasi Euler, dan galatnya menyusut separuh setiap panjang
langkah diperkecil separuh.
\end{kotak}
\end{frame}

% ========================================================= bagian 6
\section{Yang Belum Ada pada Fisika Ini}
\sectionframe{06}{Yang Belum Ada pada Fisika Ini}
  {Batas kemampuan model dan pengait ke materi berikutnya}

\begin{frame}{Tabrakan yang Baru Diketahui Sebagian}
\awalan{Program bab ini hanya mengenal dua jenis batas: lantai sebagai satu
garis mendatar dan kedua dinding sebagai dua garis tegak.}
\begin{itemize}
  \item Pemeriksaannya membandingkan satu nilai, bukan memeriksa bentuk.
  \item Permainan yang sesungguhnya memiliki rintangan di tengah ruang
        permainan, benda yang bertumpuk, dan benda yang bergerak saling
        mendekat.
  \item Pemeriksaan semacam itu menuntut cara yang lebih umum: membandingkan
        persegi pembatas dua benda, menentukan dari sisi mana tabrakan terjadi,
        lalu menyelesaikannya per sumbu. Seluruhnya dibahas pada Bab VII.
  \item Pekerjaan bab ini tidak sia-sia: urutan ``perbaiki posisi lebih dahulu,
        baru ubah kecepatan'' dan pemeriksaan dengan tanda ketidaksamaan tetap
        berlaku pada pemeriksaan tabrakan yang lebih umum.
\end{itemize}
\begin{kiat}
Keterbatasan itu menjadi titik tolak Materi 5: \emph{sprite}, \emph{group}, dan
\emph{collision detection} memberi cara yang berlaku untuk benda apa pun,
bukan hanya untuk tepi jendela.
\end{kiat}
\end{frame}

\begin{frame}{Sifat Benda yang Belum Dimodelkan}
\awalan{Tiga sifat berikut belum ada pada model ini, dan masing-masing menuntut
tambahan tersendiri.}
{\scriptsize
\renewcommand{\arraystretch}{1.16}
\begin{tabularx}{\textwidth}{@{}>{\raggedright\arraybackslash}p{3.0cm}
    >{\raggedright\arraybackslash}p{5.3cm}
    >{\raggedright\arraybackslash}X@{}}
\toprule
\textbf{Sifat} & \textbf{Yang seharusnya terjadi} & \textbf{Akibat pada model
  ini} \\
\midrule
Gesekan (\emph{friction}) &
  Benda yang bergeser di atas lantai melambat dan akhirnya berhenti. &
  Kecepatan mendatar hanya berubah karena tombol, sehingga pemain dapat
  meluncur selamanya. \\
\addlinespace
Hambatan udara &
  Benda yang jatuh sangat lama mencapai kecepatan akhir dan tidak bertambah
  cepat lagi. &
  Kecepatan bertambah tanpa batas; karena itu perlu batas kecepatan untuk
  mencegah benda menembus lantai. \\
\addlinespace
Permukaan miring dan bertingkat &
  Arah normal menentukan gerak, dan perpindahan dibagi menjadi beberapa langkah
  kecil. &
  Lantai hanya berupa satu garis mendatar; belum ada dukungan untuk kemiringan
  maupun tangga. \\
\bottomrule
\end{tabularx}}
\vspace{0.14cm}
\begin{kotak}
Ketiganya tidak menghalangi permainan sederhana. Namun begitu level memuat
tanjakan atau benda yang meluncur, model ini perlu dilengkapi.
\end{kotak}
\end{frame}

\begin{frame}{Perhitungan Diserahkan kepada Mesin}
\awalan{Seperti seluruh Blok I, bab ini memperlihatkan apa yang sesungguhnya
dikerjakan oleh mesin permainan. Pada Unity nanti, gravitasi, pantulan, dan
tabrakan dihitung oleh \emph{engine} dengan cara yang jauh lebih cermat.}
{\scriptsize
\renewcommand{\arraystretch}{1.16}
\begin{tabularx}{\textwidth}{@{}>{\raggedright\arraybackslash}p{2.2cm}
    >{\raggedright\arraybackslash}p{5.6cm}
    >{\raggedright\arraybackslash}X@{}}
\toprule
\textbf{Bab} & \textbf{Yang ditambahkan} & \textbf{Akibatnya} \\
\midrule
Bab VII & Deteksi tabrakan antar-objek, \emph{sprite} dan \emph{group} &
  Benda dapat bertabrakan dengan rintangan dan dengan benda lain, bukan hanya
  dengan tepi jendela. \\
\addlinespace
Bab VIII--XI & Pengendalian NPC &
  Kecepatan tidak lagi berasal dari tombol atau bernilai tetap, melainkan dari
  keputusan yang diambil program sendiri. \\
\addlinespace
Bab XIII & Fisika mesin pada Unity &
  Perhitungan percepatan, tabrakan, dan gesekan diserahkan kepada mesin,
  sedangkan kode Anda mengatur pengaturan dan menanggapi peristiwanya. \\
\bottomrule
\end{tabularx}}
\vspace{0.14cm}
\begin{kotak}
Yang tidak berubah adalah cara berpikirnya: benda memiliki posisi dan kecepatan,
percepatan mengubah kecepatan, kecepatan mengubah posisi, dan tabrakan
menyelesaikan posisi lebih dahulu sebelum mengubah kecepatan.
\end{kotak}
\end{frame}

% --------------------------------------------------------- aktivitas
\begin{frame}{Aktivitas di Kelas \textperiodcentered\ Praktikum 04}
\awalan{Pengalaman belajar minggu 4 pada RPS, dikerjakan berurutan di kelas lalu
diselesaikan di rumah.}
\begin{enumerate}
  \item Ketik ulang \texttt{ch06\_fisika.py} --- jangan salin-tempel ---
        jalankan, lalu amati pemain yang berhenti di lantai dan bola yang
        memantul makin pendek.
  \item Uji lompatan: pastikan pemain hanya melompat ketika menyentuh lantai,
        dan menahan tombol tidak menghasilkan lompatan berulang.
  \item Jalankan pada 30, 60, dan 120 FPS, lalu catat tinggi lompatan terukur
        dan tinggi menurut hitungan pada lembar kerja.
  \item Ubah rasa lompatan dengan tiga pasangan nilai:
        \texttt{-{}-gravity 1200 -{}-jump-speed 600},
        \texttt{-{}-gravity 1800 -{}-jump-speed 735}, dan
        \texttt{-{}-gravity 2400 -{}-jump-speed 849}.
  \item Ubah koefisien restitusi bola menjadi 0,5; 0,95; dan 1,0. Amati lama
        bola memantul pada masing-masing nilai.
  \item Tambahkan satu bongkahan sebagai lantai kedua, atur agar pemain dapat
        mendarat di atasnya, lalu tulis catatan dan simpan pekerjaan ke Git.
\end{enumerate}
\vspace{0.12cm}
{\footnotesize\color{ubAccent}\itshape Rincian lengkap: Praktikum 04 pada Bab VI
buku ajar. Kumpulkan melalui repositori Git pribadi.}
\end{frame}

% --------------------------------------------------------- rangkuman
\begin{frame}{Rangkuman}
\begin{itemize}
  \item Percepatan adalah tingkat kedua pada pembaruan: percepatan mengubah
        kecepatan, lalu kecepatan mengubah posisi. Urutan itu wajib, sesuai
        kaidah Euler semi-implisit.
  \item Pada koordinat layar, gravitasi bernilai positif karena mengarah ke
        bawah, sedangkan kecepatan lompatan bernilai negatif karena mengarah
        ke atas.
  \item Lantai diperiksa dengan tanda ketidaksamaan, bukan kesamaan; posisi
        dikembalikan ke permukaan lantai sebelum kecepatan diubah.
  \item Pemain memakai $e = 0$ agar berhenti seketika dan kendalinya andal;
        bola memakai $0 < e < 1$ agar pantulannya mereda dengan sendirinya.
  \item Lompatan adalah impuls: kecepatan diganti satu baris, dibaca sebagai
        peristiwa \texttt{KEYDOWN}, dan hanya diizinkan di tanah. Tinggi
        lompatan $h = v_0^2/(2g)$ dan lama di udara $t = 2v_0/g$.
  \item Pantulan mengikuti $v' = -e\,v$; tinggi menyusut $e^2$ kali (pada
        $e = 0{,}8$ tinggal 64 persen), pantulan kecil dihentikan, dan galat
        Euler pada 60 FPS sekitar empat persen.
\end{itemize}
\end{frame}

% --------------------------------------------------------- cek
\begin{frame}{Cek Pemahaman}
\begin{enumerate}
  \item Tuliskan dua baris integrasi Euler, lalu jelaskan mengapa urutannya
        tidak boleh dibalik.
  \item Mengapa pemeriksaan lantai memakai tanda ketidaksamaan, bukan tanda
        kesamaan? Apa yang terjadi pada benda yang bergerak cepat?
  \item Sebuah bola dijatuhkan dari ketinggian 120 piksel dengan koefisien
        restitusi $0{,}6$. Hitunglah tinggi pantulan pertamanya.
\end{enumerate}
\vspace{0.35cm}
{\footnotesize\color{ubMuted}\itshape Jawab tanpa membuka catatan, lalu
bandingkan dengan pasangan di sebelah Anda.}
\end{frame}

% --------------------------------------------------------- penutup
\penutup{Tugas Praktikum 04}
  {Pertemuan berikutnya mengorganisasi benda ke dalam \emph{sprite} dan
   \emph{group}, lalu mendeteksi tabrakan dengan \texttt{spritecollide},
   \texttt{groupcollide}, dan AABB berbasis \texttt{Rect}. Urutan ``perbaiki
   posisi lebih dahulu, baru ubah kecepatan'' tetap berlaku.}

\end{document}
